% GENERATED FILE. Edit canonical lessons, then run npm run generate:books.
% canonical-source-hash: fnv1a64:07311baa843c066d
% canonical-lessons: FR-C66-dessiner, FR-C66-peindre, FR-C66-compter, FR-C66-gagner, FR-C66-cacher

\chapter{Draw, Paint, Count, Win, Hide}
\label{ch:fr-draw-paint-count-win-hide}
\hlchaptermodality{66}

\begin{chapteropening}
\textbf{By the end of this chapter:} \emph{I can say to draw, to paint, to count, to win, to hide.}

Complete the last lesson of chapter 66: I can say to draw, to paint, to count, to win, to hide.
\end{chapteropening}

\section[dessiner]{dessiner — to draw}

\label{lesson:FR-C66-dessiner}



\begin{quote}
Before the new one: say the French for to receive, then the French for to feel; to smell.
\end{quote}

\subsection*{You'll want to know: dessiner}
\textbf{dessiner} — \textquotedblleft{}to draw\textquotedblright{}.

\begin{cousinweb}
\textbf{dessiner} $\leftarrow$ Middle French \emph{desseigner}, from Italian \emph{disegnare}, \textquotedblleft{}to draw\textquotedblright{}, $\leftarrow$ Latin \emph{designare}, \textquotedblleft{}to mark out\textquotedblright{} — borrowed from Italian. (EtymDB 2.1 records the Middle French and Latin forms.)
\end{cousinweb}

\begin{grammarlens}[title={What you've built}]
One more word for this chapter.
\end{grammarlens}

\subsection*{Your turn}
\begin{itemize}
\raggedright
  \item \emph{Say it:} \emph{dessiner}
  \item \emph{Say it:} \emph{dessiner}, once more
  \item \emph{Say it:} say \emph{dessiner}
  \item \emph{Recall:} say the French for to receive, then the French for to feel; to smell, then say \emph{dessiner} again
\end{itemize}

\begin{tcolorbox}[breakable,colback=teal!4,colframe=teal!35!black,title={Before you move on}]
What does dessiner mean? (\textquotedblleft{}To draw\textquotedblright{}.) Say it once more.
\end{tcolorbox}
\section[peindre]{peindre — to paint}

\label{lesson:FR-C66-peindre}



\begin{quote}
Before the new one: say the French for to feel; to smell, then the French for to draw.
\end{quote}

\subsection*{You'll want to know: peindre}
\textbf{peindre} — \textquotedblleft{}to paint\textquotedblright{}.

\begin{cousinweb}
\textbf{peindre} $\leftarrow$ Old French \emph{peindre} $\leftarrow$ Latin \emph{pingere}, \textquotedblleft{}to paint\textquotedblright{} — inherited. (EtymDB 2.1)
\end{cousinweb}

\begin{grammarlens}[title={What you've built}]
One more word for this chapter.
\end{grammarlens}

\subsection*{Your turn}
\begin{itemize}
\raggedright
  \item \emph{Say it:} \emph{peindre}
  \item \emph{Say it:} \emph{peindre}, once more
  \item \emph{Say it:} say \emph{peindre}
  \item \emph{Recall:} say the French for to feel; to smell, then the French for to draw, then say \emph{peindre} again
\end{itemize}

\begin{tcolorbox}[breakable,colback=teal!4,colframe=teal!35!black,title={Before you move on}]
What does peindre mean? (\textquotedblleft{}To paint\textquotedblright{}.) Say it once more.
\end{tcolorbox}
\section[compter]{compter — to count}

\label{lesson:FR-C66-compter}



\begin{quote}
Before the new one: say the French for to draw, then the French for to paint.
\end{quote}

\subsection*{You'll want to know: compter}
\textbf{compter} — \textquotedblleft{}to count\textquotedblright{}.

\begin{cousinweb}
\textbf{compter} $\leftarrow$ Latin \emph{computare}, \textquotedblleft{}to reckon\textquotedblright{} — inherited. (EtymDB 2.1) Old French wrote it \emph{conter}; the \emph{mp} was put back later to recall the Latin.
\end{cousinweb}

\begin{grammarlens}[title={What you've built}]
One more word for this chapter.
\end{grammarlens}

\subsection*{Your turn}
\begin{itemize}
\raggedright
  \item \emph{Say it:} \emph{compter}
  \item \emph{Say it:} \emph{compter}, once more
  \item \emph{Say it:} say \emph{compter}
  \item \emph{Recall:} say the French for to draw, then the French for to paint, then say \emph{compter} again
\end{itemize}

\begin{tcolorbox}[breakable,colback=teal!4,colframe=teal!35!black,title={Before you move on}]
What does compter mean? (\textquotedblleft{}To count\textquotedblright{}.) Say it once more.
\end{tcolorbox}
\section[gagner]{gagner — to win; to earn}

\label{lesson:FR-C66-gagner}



\begin{quote}
Before the new one: say the French for to paint, then the French for to count.
\end{quote}

\subsection*{You'll want to know: gagner}
\textbf{gagner} — \textquotedblleft{}to win; to earn\textquotedblright{}.

\begin{cousinweb}
\textbf{gagner} $\leftarrow$ Old French \emph{gaaignier} $\leftarrow$ Frankish \emph{*waidanjan}, \textquotedblleft{}to hunt for food, to graze\textquotedblright{} — borrowed from the Franks. (EtymDB 2.1)
\end{cousinweb}

\begin{grammarlens}[title={What you've built}]
One more word for this chapter.
\end{grammarlens}

\subsection*{Your turn}
\begin{itemize}
\raggedright
  \item \emph{Say it:} \emph{gagner}
  \item \emph{Say it:} \emph{gagner}, once more
  \item \emph{Say it:} say \emph{gagner}
  \item \emph{Recall:} say the French for to paint, then the French for to count, then say \emph{gagner} again
\end{itemize}

\begin{tcolorbox}[breakable,colback=teal!4,colframe=teal!35!black,title={Before you move on}]
What does gagner mean? (\textquotedblleft{}To win; to earn\textquotedblright{}.) Say it once more.
\end{tcolorbox}
\section[cacher]{cacher — to hide}

\label{lesson:FR-C66-cacher}



\begin{quote}
Before the new one: say the French for to count, then the French for to win; to earn.
\end{quote}

\subsection*{You'll want to know: cacher}
\textbf{cacher} — \textquotedblleft{}to hide\textquotedblright{}.

\begin{cousinweb}
\textbf{cacher} $\leftarrow$ Old French \emph{cachier} $\leftarrow$ Vulgar Latin \emph{*coacticare}, \textquotedblleft{}to press together\textquotedblright{} — inherited. (EtymDB 2.1)
\end{cousinweb}

\begin{grammarlens}[title={What you've built}]
One more word for this chapter.
\end{grammarlens}

\subsection*{Your turn}
\begin{itemize}
\raggedright
  \item \emph{Say it:} \emph{cacher}
  \item \emph{Say it:} \emph{cacher}, once more
  \item \emph{Say it:} say \emph{cacher}
  \item \emph{Recall:} say the French for to count, then the French for to win; to earn, then say \emph{cacher} again
\end{itemize}

\begin{tcolorbox}[breakable,colback=teal!4,colframe=teal!35!black,title={Before you move on}]
What does cacher mean? (\textquotedblleft{}To hide\textquotedblright{}.) Say it once more.
\end{tcolorbox}
